\documentclass{scrbook} 

\usepackage{xltxtra,tabularx}

\begin{document}

\subsection*{Mäˊrjj lij puõcci}
Ko Evvan puätt pâˊsspeeiˊv Määˊrj årra, son jordd, što nijdd reäkk. Suu čââˊlm da njuuˊnn âˊvvee da lie sami ruõpssâd.\\

\begin{tabularx}{\textwidth}{l X}
Evvan:& 	Mii tuˊst lij?\\
Mäˊrjj:& 	Jiõm tieˊđ. Muˊst čââˊlm da njuuˊnn âˊvvee da vueiˊvv čuâcc.\\
Evvan:&	Lij-a tuˊst žaar?\\
Mäˊrjj:&	Vuâsti lij, ko mon leäm čuuˊt viõǥǥteˊm da kâlmstõõlam.\\
Evvan:&	Leäk-a soiˊttjam dåhttra?\\
Mäˊrjj:&	Jiõm leäkku.\\
Evvan:&	Ââˊn âlgg soiˊttjed. Mon ooʒʒam nââmar da ton suåvak ääiˊj.\\
\end{tabularx}\\

Evvan soiˊttai tiõrvâsvuõttpõˊrtte, leâˊša Mäˊrjj lij nuˊtt puõcci, što son ij vueiˊt mainsted. Evvan maainast suu peäˊlest dåhttrin. Tät päkk pueˊtted tâˊlles tiõrvâsvuõttpõˊrtte. Suäna mâˊnne Evvan aautin. Dåhttar tuˊtǩǩad Määˊrj da ceälkk: “Koozz čuâcc? Lij-a tuˊst nuõpp, ko njuuˊnn nuˊtt âˊvvad?”

Mäˊrjj ceälkk, što suˊst lij čåuˊjjkõpp da suu loojjat. Njuuˊnn âˊvvad miâlggâd puk ääiˊj. Suˊst sätt leeˊd žaar še. Dåhttar kââčč, što lie-a Määˊrjest šâˊlddrään, lie-a suˊst kaass leˊbe piânnai, lij-a põõrtâst jiânnai põõˊl. Son ceälkk, što Määˊrjest sätt leeˊd alleerǥ. Mäˊrjj ij ååsk tõn, leâˊša dåhttar tätt väˊldded suˊst võõr. Son oudd talkksid da ceälkk: “Muuˊšt väˊldded talkkâz tueˊlää da jeäˊǩǩää. De neäˊttel mâŋŋa puäđak eˊpet tiiˊǩ, što ǩiõččâp mâiˊd võõrâst kaunnâp.”

Evvan tätt pueˊtted hoiˊddjed Määˊrj. Son puätt jieˊnnines da vueˊbbines Määˊrj årra. Sij pâˊsse šââˊldid da ǩeˊtte veär. Leâˊša Mäˊrjj ij vueiˊt poorrâd.\\

\begin{tabularx}{\textwidth}{l X}
Mäˊrjj:& 	Jiõm vueiˊt poorrâd ââˊn. Jâđđa mon vuâsti vuäitam.\\
Evvan jeäˊnn:&	Mij pueˊttep eˊpet jâđđa. Âlgg-a puˊhtted mâiˊd-ne?\\
Mäˊrjj:&	Späˊsseb, jiõm taarbâž ni mâiˊd. Evvan, puäˊđ tuˊvdded muu, ko čuuˊt lij kõõumâs.\\
\end{tabularx}\\

Evvan oudd suˊnne puõˊlli mueˊrjjsääˊpp da toudd suu puârast.

\end{document}